\section{Introduction}
\label{sec:intro}

When an agent can consult multiple information sources before making a decision, a natural question is: does consulting one source make another more or less valuable?
Sources that make each other more valuable are \emph{complements}; sources that diminish each other's value are \emph{substitutes}.

\citet{ChenWaggoner2017} formalized this distinction via the pointwise interaction $\Delta\VoI(j \mid i, b)$: the change in channel~$j$'s value of information at belief~$b$ after observing channel~$i$.
Their work established global conditions (across all beliefs) for substitution and complementarity. Our contribution is to localize these conditions: we characterize \emph{which beliefs} give rise to complementarity and which to substitution, for a fixed pair of channels and a fixed decision problem.
Concretely, we establish three results:
\begin{enumerate}
  \item \textbf{Bregman decomposition} (Proposition~\ref{prop:decomposition}): $\Delta\VoI(j \mid i, b)$ decomposes into a non-negative complement force $\E[D_g(\bpost, b)]$ minus a non-negative substitute force $\E[D_h(\bpost, b)]$, where $D_g$ and $D_h$ are Bregman divergences of auxiliary convex functions.
  \item \textbf{Interior complementarity} (Theorem~\ref{thm:interior}): If all of channel~$i$'s posteriors lie within the current decision region, then $\Delta\VoI(j \mid i, b) \geq 0$.
  \item \textbf{Converse} (Theorem~\ref{thm:converse}): If $\Delta\VoI(j \mid i, b) < 0$, at least one posterior of channel~$i$ lies in a different decision region from~$b$.
\end{enumerate}

Boundary crossing is necessary for substitution but not sufficient: channels can cross decision boundaries and still complement each other. Section~\ref{sec:setup} formalizes the setting; Section~\ref{sec:localization} presents the main results; Section~\ref{sec:example} illustrates them on a concrete example.

\begin{figure}[htbp]
\centering
\includegraphics[width=\textwidth]{figures/interaction_decomposition.pdf}
\caption{Belief-simplex geometry of the interaction.
\textit{Left}: the sign of $\Delta\VoI$ partitions the simplex into complement and substitute regions.
\textit{Middle, Right}: the two constituent forces; the complement force is diffuse across the interior while the substitute force concentrates sharply along decision boundaries. Their difference determines the sign of $\Delta\VoI$.}
\label{fig:hero}
\end{figure}
